\textbf{Phần 1. Trắc nghiệm lựa chọn}
\setcounterpageref{ex}{0}

\begin{ex} %[5D4N22]
Kết quả của phép tính $0{,}36 : 9$ là
\choice
{\True $0{,}04$}
{$0{,}4$}
{$4$}
{$0{,}004$}
\loigiai{
Ta có: $0{,}36 : 9 = 0{,}04$.\\
\textbf{Đáp án của câu hỏi là A}
}
\end{ex}

\begin{ex} %[5D4N22]
Kết quả của phép tính $26 : 8$ là
\choice
{$3{,}2$}
{\True $3{,}25$}
{$3{,}5$}
{$3{,}125$}
\loigiai{
Thực hiện chia $26 : 8 = 3{,}25$.\\
\textbf{Đáp án của câu hỏi là B}
}
\end{ex}

\begin{ex} %[5D4H22]
Kết quả của phép tính $19{,}95 : 19$ là
\choice
{$1{,}5$}
{$1{,}005$}
{$10{,}5$}
{\True $1{,}05$}
\loigiai{
Ta có: $19 : 19 = 1$; $9 : 19 = 0$; $95 : 19 = 5$. Kết quả là $1{,}05$.\\
\textbf{Đáp án của câu hỏi là D}
}
\end{ex}

\begin{ex} %[5D4H22]
Kết quả của phép tính $57 : 9{,}5$ là
\choice
{\True $6$}
{$0{,}6$}
{$60$}
{$600$}
\loigiai{
Chuyển về chia hai số tự nhiên: $570 : 95 = 6$.\\
\textbf{Đáp án của câu hỏi là A}
}
\end{ex}

\begin{ex} %[5D4H22]
Kết quả của phép tính $128 : 0{,}25$ là
\choice
{$32$}
{$51{,}2$}
{\True $512$}
{$5\,120$}
\loigiai{
Ta có: $128 : 0{,}25 = 12\,800 : 25 = 512$ (hoặc $128 \times 4 = 512$).\\
\textbf{Đáp án của câu hỏi là C}
}
\end{ex}

\begin{ex} %[5D4H22]
Kết quả của phép tính $48{,}3 : 3{,}5$ là
\choice
{$13{,}5$}
{\True $13{,}8$}
{$138$}
{$1{,}38$}
\loigiai{
Ta có: $48{,}3 : 3{,}5 = 483 : 35 = 13{,}8$.\\
\textbf{Đáp án của câu hỏi là B}
}
\end{ex}

\begin{ex} %[5D4H22]
Cho biết $7\,657 : 31 = 247$. Không thực hiện phép tính, giá trị của $76{,}57 : 31$ là
\choice
{$24{,}7$}
{$0{,}247$}
{\True $2{,}47$}
{$247$}
\loigiai{
Số bị chia giảm đi $100$ lần thì thương cũng giảm đi $100$ lần: $247 : 100 = 2{,}47$.\\
\textbf{Đáp án của câu hỏi là C}
}
\end{ex}

\begin{ex} %[5D4H22]
Một con tàu vũ trụ trong $2$ giây bay được $103\text{ km}$. Vận tốc trung bình của con tàu là
\choice
{$50{,}5\text{ km/giây}$}
{$52{,}5\text{ km/giây}$}
{$206\text{ km/giây}$}
{\True $51{,}5\text{ km/giây}$}
\loigiai{
Vận tốc trung bình của tàu là: $103 : 2 = 51{,}5\text{ (km/giây)}$.\\
\textbf{Đáp án của câu hỏi là D}
}
\end{ex}

\begin{ex} %[5D4V22]
Một nhà khảo cổ đã đổ hết $15\text{ l}$ dầu vào các rô-bốt chuột chũi, mỗi rô-bốt chứa được $0{,}75\text{ l}$ dầu. Hỏi có tất cả bao nhiêu rô-bốt chuột chũi?
\choice
{\True $20$ rô-bốt}
{$15$ rô-bốt}
{$25$ rô-bốt}
{$30$ rô-bốt}
\loigiai{
Số rô-bốt chuột chũi là: $15 : 0{,}75 = 1\,500 : 75 = 20$ (rô-bốt).\\
\textbf{Đáp án của câu hỏi là A}
}
\end{ex}

\begin{ex} %[5D4V22]
Trên $1{,}5\text{ ha}$ đất người ta thu hoạch được $3\text{ tấn}$ hạt điều thô. Hỏi trung bình trên mỗi héc-ta đất đó thu hoạch được bao nhiêu tấn hạt điều thô?
\choice
{$1{,}5\text{ tấn}$}
{\True $2\text{ tấn}$}
{$4{,}5\text{ tấn}$}
{$0{,}5\text{ tấn}$}
\loigiai{
Trên mỗi héc-ta thu hoạch được: $3 : 1{,}5 = 2\text{ (tấn)}$.\\
\textbf{Đáp án của câu hỏi là B}
}
\end{ex}

\begin{ex} %[5D4V22]
Một chú rồng nhỏ $4$ chiếc răng sâu và phải trả cho nha sĩ tất cả $15{,}4\text{ kg}$ kẹo. Biết số kẹo trả cho mỗi chiếc răng là như nhau. Hỏi nhổ mỗi chiếc răng sâu chú rồng phải trả bao nhiêu ki-lô-gam kẹo?
\choice
{$3{,}6\text{ kg}$}
{$3{,}8\text{ kg}$}
{$3{,}95\text{ kg}$}
{\True $3{,}85\text{ kg}$}
\loigiai{
Số kẹo cho mỗi chiếc răng sâu là: $15{,}4 : 4 = 3{,}85\text{ (kg)}$.\\
\textbf{Đáp án của câu hỏi là D}
}
\end{ex}

\begin{ex} %[5D4V22]
Mặt sàn một nhà kính trồng rau dạng hình chữ nhật có chiều dài $51{,}5\text{ m}$ và chiều rộng bằng $\dfrac{2}{5}$ chiều dài. Chu vi mặt sàn nhà kính đó là
\choice
{$72{,}1\text{ m}$}
{$103\text{ m}$}
{\True $144{,}2\text{ m}$}
{$144\text{ m}$}
\loigiai{
Chiều rộng mặt sàn là: $51{,}5 \times \dfrac{2}{5} = 20{,}6\text{ (m)}$.\\
Chu vi mặt sàn là: $(51{,}5 + 20{,}6) \times 2 = 72{,}1 \times 2 = 144{,}2\text{ (m)}$.\\
\textbf{Đáp án của câu hỏi là C}
}
\end{ex}

\textbf{Phần 2. Trắc nghiệm đúng sai}
\setcounterpageref{ex}{0}

\begin{ex} %[5D4H22]
Xét tính đúng sai của các phép chia sau:
\choiceTF
{\True Phép chia $15 : 8 = 1{,}875$.}
{Phép chia $6 : 25 = 2{,}4$.}
{\True Phép chia $702 : 7{,}2 = 97{,}5$.}
{\True Phép chia $5{,}4 : 0{,}25 = 21{,}6$.}
\loigiai{
\begin{itemize}
    \item $15 : 8 = 1{,}875$.
    \item $6 : 25 = 0{,}24 \neq 2{,}4$.
    \item $702 : 7{,}2 = 7\,020 : 72 = 97{,}5$.
    \item $5{,}4 : 0{,}25 = 540 : 25 = 21{,}6$.
\end{itemize}
\begin{itemize}
    \item a) \textbf{Đúng}.
    \item b) \textbf{Sai}: $6 : 25 = 0{,}24$.
    \item c) \textbf{Đúng}.
    \item d) \textbf{Đúng}.
\end{itemize}
}
\end{ex}

\begin{ex} %[5D4V22]
Rô-bốt chia đều $9{,}68\text{ yến}$ cá vào $8$ khay cá bằng nhau. Xét tính đúng sai của các khẳng định sau:
\choiceTF
{\True Mỗi khay cá đựng $1{,}21\text{ yến}$ cá.}
{Đổi $1{,}21\text{ yến} = 121\text{ kg}$.}
{\True Khối lượng cá đựng trong $5$ khay là $6{,}05\text{ yến}$.}
{Cần dùng ít nhất $10$ khay nếu mỗi khay chỉ đựng đúng $1\text{ yến}$ cá.}
\loigiai{
\begin{itemize}
    \item Khối lượng mỗi khay: $9{,}68 : 8 = 1{,}21\text{ yến}$.
    \item $1{,}21\text{ yến} = 1{,}21 \times 10 = 12{,}1\text{ kg} \neq 121\text{ kg}$.
    \item $5$ khay đựng: $5 \times 1{,}21 = 6{,}05\text{ yến}$.
    \item $9{,}68 : 1 = 9{,}68 \Rightarrow$ cần $10$ khay (đúng, vì $9$ khay chỉ chứa được $9\text{ yến} < 9{,}68\text{ yến}$).
\end{itemize}
\begin{itemize}
    \item a) \textbf{Đúng}.
    \item b) \textbf{Sai}: $1{,}21\text{ yến} = 12{,}1\text{ kg}$.
    \item c) \textbf{Đúng}.
    \item d) \textbf{Đúng}.
\end{itemize}
}
\end{ex}

\textbf{Phần 3. Trả lời ngắn}
\setcounterpageref{ex}{0}

\begin{ex} %[5D4H22]
Tính thương của phép chia: $95{,}2 : 68$. (Nhập kết quả dưới dạng số thập phân).
\shortans{1,4}
\loigiai{
\textbf{Đáp số:} $1{,}4$\\
$95{,}2 : 68 = 1{,}4$.
}
\end{ex}

\begin{ex} %[5D4H22]
Tính thương của phép chia: $54 : 0{,}45$.
\shortans{120}
\loigiai{
\textbf{Đáp số:} $120$\\
$54 : 0{,}45 = 5\,400 : 45 = 120$.
}
\end{ex}

\begin{ex} %[5D4H22]
Tính thương của phép chia: $144{,}55 : 3{,}5$. (Nhập kết quả dưới dạng số thập phân).
\shortans{41,3}
\loigiai{
\textbf{Đáp số:} $41{,}3$\\
$144{,}55 : 3{,}5 = 1\,445{,}5 : 35 = 41{,}3$.
}
\end{ex}

\begin{ex} %[5D4V22]
Một thùng có $15\text{ l}$ dầu. Người ta rót vào các chai nhỏ, mỗi chai chứa $0{,}75\text{ l}$. Hỏi rót được tất cả bao nhiêu chai dầu?
\shortans{20}
\loigiai{
\textbf{Đáp số:} $20$\\
$15 : 0{,}75 = 20$ (chai).
}
\end{ex}

\begin{ex} %[5D4V22]
Nhổ $4$ chiếc răng sâu hết $15{,}4\text{ kg}$ kẹo. Hỏi nhổ $1$ chiếc răng sâu như vậy hết bao nhiêu ki-lô-gam kẹo? (Nhập kết quả dưới dạng số thập phân).
\shortans{3,85}
\loigiai{
\textbf{Đáp số:} $3{,}85$\\
$15{,}4 : 4 = 3{,}85\text{ (kg)}$.
}
\end{ex}

\begin{ex} %[5D4V22]
Mặt sàn hình chữ nhật dài $51{,}5\text{ m}$, rộng bằng $\dfrac{2}{5}$ chiều dài. Tính chu vi mặt sàn đó theo mét. (Nhập kết quả dưới dạng số thập phân).
\shortans{144,2}
\loigiai{
\textbf{Đáp số:} $144{,}2$\\
Rộng: $51{,}5 \times \dfrac{2}{5} = 20{,}6\text{ m}$. Chu vi: $(51{,}5 + 20{,}6) \times 2 = 144{,}2\text{ m}$.
}
\end{ex}
